% !TEX spellcheck = en_US
\documentclass[aspectratio=169,10pt,compress]{beamer}
\usepackage{../../cis000_latex/lib/hblibPresentation}
\usepackage{../../cis000_latex/lib/hblibCodeoutput}
\usepackage{../../cis000_latex/lib/hblibDatastructures}
\usepackage[T1]{fontenc}
\usepackage{lmodern}
\usepackage{booktabs}

\lstset{columns=fixed, breaklines=true, breakatwhitespace=false}

\title{Linked Stacks, Queues, and Deques}
\subtitle{Sophomore Level: Generic Containers and Doubly-Linked Deque}
\author[Bingol]{Ruhan Alpaydin}
\institute[FCIS]{Faculty of Computer and Information Sciences\\Yeditepe University}
\date{{\tiny \hbTimeStamp}}

\begin{document}

\begin{frame}[t,fragile]\titlepage\end{frame}

\begin{frame}[t,fragile]\frametitle{Outline}
  \begin{columns}[T]
    \begin{column}{0.48\textwidth}\tableofcontents[sections={1-3},hideallsubsections]\end{column}
    \begin{column}{0.48\textwidth}\tableofcontents[sections={4-5},hideallsubsections]\end{column}
  \end{columns}
\end{frame}

\hSection[ADT]{The Interfaces}
\begin{frame}[t,fragile]\frametitle{\texttt{StackADT<T>} and \texttt{QueueADT<T>}}
  \begin{columns}[T]
    \begin{column}{0.48\textwidth}
      \lstcode[basicstyle=\ttfamily\scriptsize]{codeoutput/StackADT.java}{Java}
    \end{column}
    \begin{column}{0.48\textwidth}
      \lstcode[basicstyle=\ttfamily\scriptsize]{codeoutput/QueueADT.java}{Java}
    \end{column}
  \end{columns}
  \vspace{0.4em}
  \small
  \hRemark{Generic interfaces} mean the same code works for \hCode{Integer}, \hCode{String}, or any custom object.
\end{frame}

\begin{frame}[t,fragile]\frametitle{\texttt{DequeADT<T>}}
  \lstcode[basicstyle=\ttfamily\scriptsize]{codeoutput/DequeADT.java}{Java}
  \begin{itemize}\small
    \item Double-Ended Queue: inserts and removes at both ends.
    \item Implements both \hCode{StackADT} and \hCode{QueueADT} behaviour.
  \end{itemize}
\end{frame}

\hSection[Deque]{Doubly Linked Deque}
\begin{frame}[t,fragile]\frametitle{The Node}
  \begin{columns}[T]
    \begin{column}{0.5\textwidth}
      \begin{lstlisting}[language=Java, basicstyle=\ttfamily\scriptsize, frame=single]
private static class Node<T> {
    T data;
    Node<T> prev, next;
    Node(T data) {
        this.data = data;
    }
}
      \end{lstlisting}
    \end{column}
    \begin{column}{0.46\textwidth}
      \begin{itemize}\small
        \item \hCode{prev} points to the previous node.
        \item \hCode{next} points to the next node.
        \item This symmetry is what gives $O(1)$ at \emph{both} ends.
      \end{itemize}
    \end{column}
  \end{columns}
\end{frame}

\begin{frame}[t,fragile]\frametitle{\texttt{addFirst} and \texttt{addLast}}
  \begin{columns}[T]
    \begin{column}{0.52\textwidth}
      \lstcode[basicstyle=\ttfamily\scriptsize, firstline=9, lastline=22, firstnumber=9]{codeoutput/LinkedDeque.java}{Java}
    \end{column}
    \begin{column}{0.44\textwidth}
      \begin{itemize}\small
        \item Both $O(1)$.
        \item Empty list: \hCode{head = tail = n}.
        \item \hCode{addFirst}: \hCode{head.prev = n; head = n}.
        \item \hCode{addLast}: \hCode{tail.next = n; tail = n}.
      \end{itemize}
    \end{column}
  \end{columns}
\end{frame}

\begin{frame}[t,fragile]\frametitle{\texttt{removeFirst} and \texttt{removeLast}}
  \begin{columns}[T]
    \begin{column}{0.52\textwidth}
      \lstcode[basicstyle=\ttfamily\scriptsize, firstline=24, lastline=39, firstnumber=24]{codeoutput/LinkedDeque.java}{Java}
    \end{column}
    \begin{column}{0.44\textwidth}
      \begin{itemize}\small
        \item \hCode{removeFirst}: move \hCode{head} right, clear \hCode{head.prev}.
        \item \hCode{removeLast}: move \hCode{tail} left, clear \hCode{tail.next}.
        \item Both $O(1)$.
      \end{itemize}
    \end{column}
  \end{columns}
\end{frame}

\begin{frame}[t,fragile]\frametitle{The run}
  \begin{columns}[T]
    \begin{column}{0.52\textwidth}
      \lstcode[basicstyle=\ttfamily\tiny]{codeoutput/LinkedDeque.java}{Java}
    \end{column}
    \begin{column}{0.44\textwidth}
      \lstterminal[basicstyle=\ttfamily\scriptsize]{codeoutput/LinkedDeque.txt}
      \small
      C is removed from the back, A and B from the front.
    \end{column}
  \end{columns}
\end{frame}

\hSection[Compare]{Comparison with JDK}
\begin{frame}[t,fragile]\frametitle{JDK equivalents}
  \begin{center}\small
  \begin{tabular}{@{}lll@{}}
    \toprule
    \textbf{Your class} & \textbf{JDK Class} & \textbf{Note} \\
    \midrule
    \hCode{LinkedDeque<T>} & \hCode{LinkedList<T>} & Implements \hCode{Deque}, \hCode{List}, \hCode{Queue} \\
    \emph{(none)} & \hCode{ArrayDeque<T>} & Faster, array-backed, no per-node overhead \\
    \bottomrule
  \end{tabular}
  \end{center}
  \vspace{0.3em}
  \begin{itemize}\small
    \item In practice, \hCode{ArrayDeque} is usually the better choice.
    \item \hCode{LinkedList} wins only when you need to insert/remove in the middle via an \hCode{Iterator}.
  \end{itemize}
\end{frame}

\begin{frame}[t,fragile]
  \begin{center}\vspace{4cm}\hRemark{\Huge Questions?}\end{center}
\end{frame}

\end{document}
